\section{Experiments}\label{experiments}

Here, we first describe the datasets (Section \ref{subsec:data}) followed by details regarding implementation and hyperparameters (Section \ref{implement}), results of our experimentations in Section \ref{sec:resl}, explainability of the proposed approach (Section \ref{explain}), and lastly a discussion of the results of our experiments (Section \ref{discuss}).

\subsection{Dataset}\label{subsec:data}


The experiments are conducted on two widely used benchmark datasets: NELL-995 \cite{Xiong2017} and FB15K-237 \cite{Toutanova2015}. Table \ref{table:KGstats} describes the specifics pretraining to the two datasets. NELL-995 is generated from the 995th iteration of the NELL system \cite{Carlson2010} by selecting the triples with the top 200 most frequent relations. FB15K-237's facts are created from FB15K \cite{Bordes2013} with redundant relations removed. Both datasets contain inverse triples to aid with path finding, allowing the reasoning agent to step backward in the graph. %All the triples with $r_q$ or $r_q^{-1}$ are removed from KG and divided into train and test samples for each reasoning task with a query relation $r_q$. 
We chose three reasoning tasks per dataset: \textit{origin}, \textit{tvLanguage}, \textit{nationality} for FB15K-237, and \textit{athletePlaysInLeague}, \textit{athletePlaysForTeam}, \textit{worksFor} for NELL-995.
%The KGs, train, and test sets are taken from \cite{Xiong2017}.
To evaluate the agent's performance, each reasoning task contains approximately ten negative examples for each positive sample, and the negative samples were generated by replacing a true claim's tail with a faked entity.

\begin{table}[h!]
\centering
\caption{Statistics of the datasets under consideration; \#Ent, \#Rel, and \#Fact denotes the number of unique entities, relations, and triples respectively}\label{Tbl:datatable}
\begin{tabular}{|l|c|c|c|}
\hline 
Dataset   & \#Ent & \#Rel & \#Fact  \\ \hline
FB15k-237 & 14,505                        & 237   & 272,115 \\ \hline
NELL-995  & 75,492                        & 200   & 154,213 \\ \hline
\end{tabular}
\label{table:KGstats}
\end{table}
%Table \ref{table:KGstats}: Statistics of the datasets. \#Ent denotes the number of unique entities and \#Rel denotes the number of relations.



\subsection{Implementation Details and Hyperparameters} \label{implement}

\noindent \textbf{Knowledge Graph Embeddings: }We use \textit{ComplEx} \cite{Trouillon2016} embedding model from \textit{Ampligraph} \cite{ampligraph} library for KG embeddings. The model's hyperparameters are chosen to balance model complexity and training time. We set the batch size as 50 and train the model for 1000 epochs. The embedding dimension is set to 20. We use Adam optimizer with a learning rate of 1e-4 and negative log-likelihood loss function. %multiclass\_nll \textcolor{red}{What is this loss function? Kindly confirm that this is correct}
%\textcolor{olive}{Correct!}. 
To prevent overfitting, we use an L3 regularization term with a regularization strength of 1e-5. We also set the random seed to 0 to ensure the reproducibility of the results. %\textcolor{red}{@Gustav There is no 80/20 training test split considered here, right? Like 80:20 split for policy network training, did we did same for KG embeddings?}
%\textcolor{olive}{I'm not sure if I will answer your question with this, but: training the KG, I used a train 'test split no unseen', where the test set contains triples which are also included in the train set. I kept it at the default value, which is 100 test triples. So the answer would be in some sense, yes, no training test split considered.}


\noindent \textbf{Policy Network: }We use REINFORCE algorithm to train the policy network and \textit{PyTorch} library for the implementation. It comprises of ReLU as an activation function in the intermediate layers and Softmax in the output layer to generate a probability distribution over the actions given the current state. The network has hidden state size of 128 neurons and was initialized using the Xavier initialization method. %\cite{Glorot2010}. 
We train the network for 100,000 episodes %\textcolor{red}{Does episodes is same as epoch?} \textcolor{olive}{No, episode is RL term. e.g one single game of tic-tac-toe at the end of which a winner is determined and awarded}
for each reasoning task, each consisting of 3 steps. We use  \textit{Adam} %\cite{Kingma2014} 
optimizer with a learning rate of 0.001. To encourage exploration, we randomly chose one of the top three most likely actions as the action to take in each state during training. \par

During the evaluation, we use beam search to obtain multiple likely actions according to the policy network's probability distribution for each state. This allows us to explore multiple possible paths through the action space and obtain a more robust estimate of the policy's performance. Three beam numbers, 3, 5, and 10 %\textcolor{red}{3, 5, and 10 are these same with the hits@k values mentioned below?}
%\textcolor{olive}{Yes, the numbers for beam numbers are the same as for hits@k}
are considered. We assess the performance of our model using an 80:20 train-test split. %\textcolor{red}{Splitting in the dataset section and also here? As in dataset section it was already mentioned we have taken the splits from a previous paper}
%\textcolor{olive}{I now think, it would be wiser to just say that we took the DATASETS from the DeepPath paper and did a 80/20 split ourselves. The double mention is because training a KG and a Policy is two different things. If you think the double mention is unneccessary, then I have no objections to removing it.}. 
We use hits@k metric, which is commonly used in knowledge graph reasoning tasks. %, with k values of 3, 5, and 10. %\textcolor{red}{3, 5, and 10 are these same with three beam numbers above?}
%\textcolor{olive}{yes}. 
 The hits@k metric measures the proportion of test samples where the correct answer is among the top k-ranked answers returned by the model. In addition, we use accuracy as a metric for the voting mechanism. This was calculated as the proportion of test samples where the weighted majority vote was the correct answer. We acknowledge that the varying degrees of information in the knowledge graphs are dependent on alternative evidential paths. However, setting thresholds to consider the number of alternative paths for each reasoning task is complex and requires further research. As such, we did not take this into account in our evaluation.

\subsection{Experimental Results}\label{sec:resl}
% Please add the following required packages to your document preamble:
% \usepackage[table,xcdraw]{xcolor}
% If you use beamer only pass "xcolor=table" option, i.e. \documentclass[xcolor=table]{beamer}

\begin{comment}
This section covers three main points:\begin{itemize}
\item Parameters considered in interpreting the results
\item Formatting of the results
\item Interpreting results for the two main tasks of the model: fact checking and explainability
\end{itemize}
\end{comment}



\noindent \textbf{Basis for interpreting the results: } The two KGs used in experimentation are different, particularly in terms of their density features, with FB15K-237 containing 18.8 triples per entity and NELL-995 containing two triples per entity. This difference may affect the model's accuracy. %, and the size of the input datasets may also influence the agent's ability to generalize knowledge. 
The subsets for both KGs consist of 3 single relation tasks, with each KG also having a combined subset containing all three relations. This diversity in task assignment may complicate the agent's task, but it seems to actually assist the agent by giving meaning to the input triple's relation. In single relation tasks, the input triple's relation never changes and therefore does not inform the agent's actions. \par

%Also, the beam number is crucial for the Beam search algorithm.% \textcolor{red}{@Gustav, will this para be more fit for next para in Results subsection? If so, kindly please move it to next para, I am finding bit difficulty here, and I may put it wrong...} \textcolor{olive}{I read both sections and I can't figure out where to move it to. Do you have a suggestion? I'll have a fresh look at this in the morning.} 
%Increasing the beam number might seem like a good way to increase the number of paths, as it  may increase the probability of reaching the target entity, but it doesn't always improve results. More paths don't necessarily lead to better accuracy, especially when searching for a single informative path.  %\textcolor{red}{@Gustav, can you refer the table where this is happening?}. \textcolor{olive}{Added table referral here:} 
%{@Gustav, can you refer the table where this is happening?}.
%Additionally, at each step of the search, all collected paths are compared to each other.


\noindent \textbf{Results: } %Table \ref{Tbl:searchfb}, Table \ref{Tbl:searchnl} and Table \ref{Tbl:votingfb}, Table \ref{Tbl:votingnl} 
Tables \ref{Tbl:searchfb} and \ref{Tbl:votingfb} show the results of our proposed approach on Hits@k and Voting Accuracy for FB15K-237. Similarly, Tables \ref{Tbl:searchnl} and \ref{Tbl:votingnl} show the results of our proposed approach on Hits@k and Voting Accuracy for NELL-995.
%shows the results of our proposed approach on Hits@k and Voting Accuracy for FB15K-237 and NELL-995 respectively. %Two accuracy metrics: search accuracy and voting accuracy %\textcolor{red}{@Gustav Is hits@k search accuracy metric, why in the result tables of Hits@k columns such as here \ref{Tbl:searchfb} it is given as acc?}.
%\textcolor{olive}{Hits@k is the search accuracy metric. In the first iteration of the paper, my thesis, I wrote search accuracy, but then I changed it to Hits@k, because it's a known metric in KG literature, and that's what the search accuracy was in essence. The mixed vocabulary is just a leftover of an unfinished rewrite.}
%\paragraph{Hits@k} the proportion of input triples for which at least one path reached the desired target entity. A higher beam number means more unique paths are considered, increasing the probability of reaching the target entity. Table II shows that considering more paths increases the proportion of input triples for which at least one path reached the target entity.
%\paragraph{Voting accuracy} is the accuracy of the proposed veracity verdict method, expressing the proportion of input triples for which the true target was predicted correctly.
%\paragraph{Beam numbers} 3, 5, and 10 are considered, with the beam number indicating how many paths are considered for one input triple.
%\subsubsection{Fact checking}
The results show that the method is able to achieve high accuracy in finding a path to the true target entity with a high enough beam number, as indicated by the results of the hits@k metric. A higher beam number means more unique paths are considered, increasing the probability of reaching the target entity. Table II shows that considering more paths increases the proportion of input triples for which at least one path reached the target entity. But that's not the case always, considering more paths can even decrease accuracy because it requires the agent to discard more paths as is the case for \textit{origin} in Table \ref{Tbl:votingfb}, and \textit{athleteplaysinleague} and \textit{worksfor} in Table \ref{Tbl:votingnl}. %\textcolor{red} %\textcolor{red}%{@Gustav, does it mean hits@k?} \textcolor{olive}{Exactly}. 
However, the accuracy drops slightly after the voting process, as the correct path must be selected from the paths found. In particular, the "origin" dataset in FB15K-237 (Table \ref{Tbl:votingfb})
\begin{comment}
\textcolor{red}{@Gustav, which table?} \textcolor{olive}{both fb15k tables, but especially voting accuracy}    
\end{comment}
 has shown to be anomalous, achieving lower accuracy than expected despite being a larger dataset. The voting task has proven to positively impact accuracy, especially for the combined dataset in FB15K-237, indicating the model's scalability. Nonetheless, the results demonstrate that voting on an answer reduces accuracy, as the correct path and target entity may be correctly identified but still be undermined by a higher majority of an entity in the vicinity of the source entity. Both FB15K-237 and NELL-995 follow similar trends in terms of results regarding beam numbers and an accuracy trade-off in voting accuracy.
\begin{comment}
\textcolor{red}{@Gustav, can you write just a couple of lines for NELL-995 dataset? Like the same trend is also observed for NELL-995 something like this.....}
\textcolor{olive}{I'm having a hard time coming up with something smart and informative to write here about NELL-995 now. Maybe in the morning... Maybe we should discuss this paragraph.} 
\end{comment}\textc

\begin{table}[h!]
\centering
\caption{Hits@k for FB15K-237; acc stands for accuracy} %; \textcolor{red}{@Gustav, why input dataset size, can we just give dataset size, and acc is accuracy?}}
\label{Tbl:searchfb}
\begin{tabular}{|c|c|c|c|c|}
\hline
            & \multicolumn{1}{l|}{dataset size} & \multicolumn{1}{l|}{acc., b=3} & \multicolumn{1}{l|}{acc., b=5} & \multicolumn{1}{l|}{acc., b=10} \\ \hline
origin      & 852                                                             & 0.464                                                      & 0.53                                                       & 0.661                                                       \\ \hline
tvLanguage  & 614                                                             & 0.52                                                       & 0.659                                                      & 0.894                                                       \\ \hline
nationality & 7134                                                            & 0.866                                                      & 0.907                                                      & 0.921                                                       \\ \hline
combined    & 8600                                                            & 0.807                                                      & 0.886                              & 0.943                               \\ \hline
\end{tabular}
\label{tabel:fb15kresults}
\end{table}

\begin{table}[h!]
\centering
\caption{Voting accuracy for FB15K-237; acc stands for accuracy} %\textcolor{red}{@Gustav, why input dataset size, can we just write dataset size, and acc is accuracy?}\textcolor{olive}{Yes and yes}}

\label{Tbl:votingfb}
\begin{tabular}
{|c|c|c|c|c|}

\hline FB15K-237   & dataset size & acc., b=3 & acc., b=5 & acc., b=10 \\ \hline
origin      & 852                & 0.357         & 0.369                                 & 0.292                                  \\ \hline
tvLanguage  & 614                & 0.463         & 0.472                                 & 0.642                                  \\ \hline
nationality & 7134               & 0.727         & 0.75                                  & 0.765                                  \\ \hline
combined    & 8600               & 0.252         & 0.602         & 0.727          \\ \hline
\end{tabular}
\end{table}

% Please add the following required packages to your document preamble:
% \usepackage[table,xcdraw]{xcolor}
% If you use beamer only pass "xcolor=table" option, i.e. \documentclass[xcolor=table]{beamer}
\begin{table}[h!]
\centering
\caption{Hits@k for NELL-995; acc stands for accuracy}
\label{Tbl:searchnl}
\resizebox{\columnwidth}{!}{\begin{tabular}
{|c|c|c|c|c|}
\hline
NELL-995     & dataset size & acc, b=3 & acc, b=5 & acc, b=10 \\ \hline
athleteplaysinleague & 12660              & 0.919         & 0.971                                 & 0.994                                  \\ \hline
athleteplayssport    & 6781               & 0.903         & 0.853                                 & 0.868                                  \\ \hline
worksfor             & 9669               & 0.748         & 0.78                                  & 0.825                                  \\ \hline
combined             & 29110              & 0.485         & 0.575         & 0.611          \\ \hline
\end{tabular}}
\end{table}

% Please add the following required packages to your document preamble:
% \usepackage[table,xcdraw]{xcolor}
% If you use beamer only pass "xcolor=table" option, i.e. \documentclass[xcolor=table]{beamer}
\begin{table}[h!]
\centering
\caption{Voting accuracy for NELL-995; acc stands for accuracy}
\label{Tbl:votingnl}
\resizebox{\columnwidth}{!}
{\begin{tabular}{|c|c|c|c|c|}
\hline
NELL-995     & dataset size & acc, b=3 & acc, b=5 & acc, b=10 \\ \hline
athleteplaysinleague & 12660              & 0.858         & 0.938                                 & 0.875                                  \\ \hline
athleteplayssport    & 6781               & 0.871         & 0.828                                 & 0.83                                   \\ \hline
worksfor             & 9669               & 0.604         & 0.59                                  & 0.509                                  \\ \hline
combined             & 29110              & 0.45          & 0.467         & 0.464          \\ \hline
\end{tabular}}
\end{table}




\subsection{Explainability}

\label{explain}

For the purpose of providing explanations, we present examples of single-relation tasks. % from different \textcolor{red}{Different why? We have only considered FB15 and NELL-995 right?} \textcolor{olive}{Correct} Knowledge Graphs (KGs), including NELL-995 Combined and FB15K-237 "origin" subsets. 
 Each example is evaluated based on the proposed approach's performance in arriving at the correct conclusion, the informativeness of the path, and the ground truth interpretation. We provide six examples in total.

%The first example shows the Agent arriving at a target entity with an informative path, but the conclusion and the presented ground truth are found to be false due to outdated information in the KG. In the second example, the Agent's reasoning is good, but the conclusion arrived at is false, indicating the need for further investigation. The third example demonstrates the Agent's ability to find a complex and informative path that leads to the true target entity. The fourth example shows the Agent being unable to move beyond the source entity, resulting in an inconclusive path. In the fifth example, the Agent finds two paths that are both plausible but only one is considered correct according to the ground truth interpretation. Finally, the sixth example shows the Agent finding two paths, one of which supports the claim indirectly but is considered incorrect according to the ground truth interpretation. By analyzing these examples, we aim to provide a clear understanding of our system's performance in different scenarios and the potential challenges in understanding and interpreting KGs.
%\newline

%\textbf{KG contains outdated information}
\noindent\textbf{Example 1: KG contains outdated information}
\\Claim: Dick Cheney works for retailstore Halliburton
\\Path: Dick Cheney – leads organization → retailstore Halliburton

\noindent The agent has found a correlation between a person leading an organization and working for the organization. This is an interesting successful case, since Dick Cheney stopped leading Halliburton the day George W. Bush took him on as vice president candidate for the 2000 US presidential elections. This means the presented ground truth is actually false due to outdated information in the KG. %\textcolor{red}{Dataset is KG, right?} \textcolor{olive}{Correct} 
\newline
%\\
%\textbf{Inconclusive path}
\noindent\textbf{Example 2: Inconclusive path}
\\Claim:	Coach John Gruden works for sportsteam Buccaneers.
\\Path:	Coach John Gruden $\leftarrow$organization hired person– sportsteam Buccaneers

\noindent The agent attempts to find a solution by considering the relation “organization hired person”.
In fact Tampa Bay Buccaneers had hired John Gruden in 2002 and Gruden worked for the team until 2008.
Thus, the claim that John Gruden works for the Buccaneers was true from 2002-2008, but is not true anymore, meaning  the agent's reasoning is good, but the conclusion arrived at is false, indicating the need for further investigation.
The agent has learned a correlation between the relations “worksfor” and “organizationhiredperson inv” meaning if an organization has hired a person, then the person works for that organization.
%\textcolor{red}{@Gustav, Do we need the next sentence or it is okay to remove it?}For this path, it might be reasonable for the agent to also look for a relation between Gruden and the Buccaneers such as “organization fired person” or “contract ended” to further check the claim.
\newline
\textbf{Example 3: Informative path with complex reasoning}
\\Claim:	Coach Brendan Shanahan plays Hockey.
\\Path:	Brendan Shanahan –plays for team→ the Devils –team plays sport→ Hockey

\noindent This is a good example of the agent finding a path that does not just contain a (near) synonym for the claimed relation (like \textit{works for} and \textit{leads organization}), but finding a more complicated path.
The agent finds that Brendan Shanahan plays for the Devils hockey team, and then finds out that the Devils play hockey.
\newline
\noindent\textbf{Example 4: Agent stays on the fence}
\\Claim:	Journalist Amir Taheri works for Los Angeles County.
\\Path:	Amir Taheri → Amir Taheri 
\newline
The Agent is unable to move beyond the source entity, resulting in an inconclusive path. This may be possible because of three reasons:
\begin{enumerate}
    \item The length of the path is fixed to three relations, no more, no less.
    \item The agent must have the possibility to decide upon a target entity and “stick with it”
    \item The agent can decide to stay put at an entity for all three transitions.
\end{enumerate}
Since most resultative paths only take one or two hops from entity to entity, the action of staying put is very common. Thus it’s always one of the top actions to choose for any state. With beam search also in place, it is very common that one path will result in the agent staying put for the entirety of the path.
\newline
\textbf{Example 5: Two true paths with only one interpreted as correct}
\\Claim: Ben Folds originates from Nashville.
\\Path 1: Ben Folds --place of birth$\rightarrow$Winston-Salem
\\Path 2: Ben Fold --location$\rightarrow$ Nashville

These paths state that Ben Folds was born in Winston-Salem and that he is located in Nashville. For this specific case, the ground truth is that Ben Folds originates from Winston-Salem, so the first path leads to the true target and the second path reinforces the "false claim".
The majority of explainable paths for the "origin" subset consist of one of two (near) synonymous relations: "place of birth" and "location".
It is also common for one of the path relations to lead to a correct target entity and the other to lead to an incorrect one, the meaning of "origin" staying quite ambiguous. The subset is concerned with musical artists, so sometimes "origin" will mean the artists place of birth, other times it will mean the city in which they became established as a musician which will often coincide with the "location" relation.
\newline
\textbf{Example 6: Two true explainable paths with only one being considered as correct}
\\Claim: Kylie Minogue originates from Australia.
\\Path 1: Kylie Minogue –place of birth→ Melbourne
\\Path 2: Kylie Minogue –place of birth→ Melbourne ←contains– Australia

\noindent The claim that Kylie Minogue originates from Australia is true in real life, but in regard to the ground truth, it is considered false.
Kylie Minogue originates from Australia which is true, but the ground truth is that Kylie Minogue originates from Melbourne, which is a city in Australia. Due to the setup of the veracity verdict mechanism, a path that concludes that Kylie Minogue originates from Australia is be considered false even though the path contains Melbourne and is very much explainable.

\subsection{Discussion}
\label{discuss}
The results of this study demonstrate that selecting an entity to be evaluated for veracity has a trade-off with accuracy. While the search algorithm may have identified the correct entity, it may not have been selected by the voting mechanism. This trade-off is not linear, and tasks with lower accuracy in identifying a correct path are more severely impacted by the results of voting. Analysis of common patterns in the paths generated by the agent highlighted several issues that should be considered in future research on this problem.

Firstly, it is crucial to use an up-to-date KG to ensure the accuracy of fact classifications. Outdated information in a KG can produce false results. Secondly, the verdict mechanism would benefit from identifying and incorporating inconclusive paths. The agent's motivation to move away from the source entity should also be increased. Additionally, paths that do not end with the target entity but support the ground truth should also be utilized to arrive at a correct conclusion. Lastly, the classification process could be improved by learning to deal with ambiguous relationships such as ``origin''. These findings suggest several avenues for future research to improve the accuracy and efficiency of veracity verdicts.